% =====================================================================
% Section 6 — Validation roadmap and future work
% =====================================================================
% Moved from Section 6 (Conclusion) on professor's recommendation:
% RESS reviewers expect a separate "Validation roadmap / Future work"
% section; the Conclusion should restate only the main findings in
% three to four sentences.
% =====================================================================

\section{Validation roadmap and future work}\label{sec:roadmap}

The deterministic indicators of Table~\ref{tbl:quant-results} and the
Monte Carlo probabilities of Table~\ref{tbl:mc-results} establish a
design-stage diagnosis but do not by themselves substitute for the
empirical confirmation of the predicted margins. This section
specifies (i)~the minimum test-rig programme that would close the
empirical gap on the AP 1.8 head, (ii)~the candidate barrier
extensions motivated by the construct-validity benchmark of
Section~\ref{sec:res-litvalidation}, and (iii)~the analytical and
optimisation directions that follow.

\subsection{Test-rig programme with statistical acceptance criteria}
\label{sec:roadmap-rig}

The minimum experimental verification programme that would close the
empirical gap on the analytical findings for the AP 1.8 head
comprises three coordinated stages. Each stage is specified with an
explicit sample size, instrumentation envelope, and statistical
acceptance rule keyed to the analytical predictions of
Section~\ref{sec:res-quant} and the probabilistic envelope of
Section~\ref{sec:res-mc}.

\begin{enumerate}
\item \textbf{Spring rig, modal verification ($n\geq 3$
specimens per spring).} A single-spring test fixture
instrumented with a non-contacting laser-displacement sensor
($\geq$ 1~kHz sampling) on the active coils identifies the surge
response across the 50--600~Hz band, under the catalogued
seat-preload boundary condition.
\emph{Acceptance criterion}: the sample-mean first surge frequency
$\bar{f}_{n}$ from $n\geq 3$ specimens shall fall within the 90\,\%
confidence interval of the analytical prediction (calibrated to the
MC distribution of Section~\ref{sec:res-mc}, i.e.\
$f_{n,\mathrm{out}}\in [304,\,334]$~Hz and
$f_{n,\mathrm{in}}\in [545,\,600]$~Hz at 90\,\% confidence). A
failure of containment is reported as a falsification of the
analytical model and triggers re-evaluation of the active-mass
hypothesis (total versus Rayleigh-effective).

\item \textbf{Motored-engine speed sweep, coil-bind and
follower-contact verification.} A motored-engine campaign from
6\,000 to 12\,000~rpm engine speed (the dynamic-loss-of-contact
boundary identified in Section~\ref{sec:res-envelope}) is
conducted with a non-contacting valve-position transducer sampled
at $\geq$~50~kHz to resolve the deceleration peak of the cam
profile. Two acceptance criteria are applied:
\emph{(2a)} absence of valve float at the deceleration peak
($a_{\max,-}=20\,362$~m\,s$^{-2}$, Section~\ref{sec:res-quant})
verified by zero-crossing analysis of the
(measured$-$predicted) lift residual at each rpm point, with a
detection threshold of $\geq$~0.05~mm departure for $\geq$~3
consecutive cam events;
\emph{(2b)} the residual coil-bind margin measured at full lift
($n\geq 5$ measurements at 10\,000~rpm) shall lie within
$\pm$0.1~mm of the MC 50\,\% credible interval
$M_{\mathrm{out}}\in [0.51,\,1.10]$~mm for the baseline outer
spring (or $M_{\mathrm{out}}\approx 2.0$~mm post-correction). A
single instance of $M_{\mathrm{out}}<0$~mm (hard contact) at any
rpm within the safe envelope falsifies the corrected design.

\item \textbf{Inner-spring fatigue campaign, Goodman criterion
verification ($n\geq 5$ specimens).} A representative sample of
inner springs is cycled under shear-stress amplitudes reproducing
the analytical
$\tau_{m,\mathrm{in}}=858$~MPa and
$\tau_{a,\mathrm{in}}=374$~MPa, with each specimen run to either
run-out at $10^{8}$ cycles or first-failure indication.
\emph{Acceptance criterion}: with $n=5$ specimens and assuming a
Weibull failure model, the upper one-sided 90\,\% confidence
bound on the empirical Goodman factor of safety
$\hat{n}_{f,\mathrm{in}}$ shall remain below unity ($\leq 0.78$
upper bound for 0 of 5 specimens reaching run-out), thereby
confirming the analytical prediction
$n_{f,\mathrm{in}}\approx 0.68$ and the MC 5\,\% quantile
$n_{f,\mathrm{in,5\%}}=0.637$. A pass condition
($\hat{n}_{f}\geq 1$ upper bound) on 5 of 5 specimens reaching
run-out would falsify the predicted Goodman deficit and require
re-evaluation of the shear endurance limit calibration.
\end{enumerate}

The total test-rig hours for this minimum programme are estimated at
the order of 600--800~h.

\subsection{Candidate barrier extensions}\label{sec:roadmap-barriers}

The construct-validity benchmark of
Section~\ref{sec:res-litvalidation} flagged two failure mechanisms
that are not currently covered by the preventive set
$P_{1}$--$P_{5}$. For each, a candidate indicator with a
literature-anchored pass criterion is proposed below so that the
framework can be expanded without re-architecting the bow diagram.

\paragraph{Soffritti, assembly/tolerance barrier (candidate
$P_{6}$).} The case of \citet{soffritti2018wornvalvetrain} attributes
premature wear at the rocker-arm/pushrod and rocker-arm/valve
interfaces to valve misalignment with respect to the seat insert.
A candidate indicator is the dimensionless \emph{stem-to-seat runout
ratio}
\begin{equation}
\rho_{\mathrm{run}}=\frac{e_{\mathrm{stem}}}{d_{\mathrm{stem}}},
\label{eq:runout}
\end{equation}
where $e_{\mathrm{stem}}$ is the measured lateral offset of the
valve-stem axis from the seat-insert axis at the assembled state
and $d_{\mathrm{stem}}$ is the valve-stem diameter. A
literature-anchored pass criterion of
$\rho_{\mathrm{run}}\leq 5\times 10^{-3}$ is proposed, consistent
with the typical $25~\mu\mathrm{m}$ alignment tolerance reported for
production valve trains \citep{heisler2002} normalised by a
nominal $d_{\mathrm{stem}}=5$--$7$~mm. The indicator is computable
from a coordinate-measuring-machine inspection of the assembled
head and can be added to Table~\ref{tbl:quant-criteria} once the
inspection protocol is documented. The barrier is not evaluated for
the AP 1.8 head in the present work but the indicator form
demonstrates that the framework is extensible to assembly-driven
threats.

\paragraph{Zhao, lubrication-regime barrier (candidate $P_{7}$).}
The case of \citet{zhao2023camshaftV6} attributes the camshaft
bearing failure to an asperity-contact regime that violates the
design assumption of full hydrodynamic lubrication. A candidate
indicator is the Stribeck-curve minimum film-thickness ratio
\begin{equation}
\Lambda_{\min}=\frac{h_{\min}}{\sqrt{R_{q,1}^{2}+R_{q,2}^{2}}},
\label{eq:stribeck}
\end{equation}
where $h_{\min}$ is the predicted minimum oil-film thickness from
the elastohydrodynamic-lubrication (EHL) calculation of the
cam--follower contact and $R_{q,1}$, $R_{q,2}$ are the root-mean-square
surface roughnesses of the two contacting surfaces. The
literature-anchored pass criterion is $\Lambda_{\min}\geq 1.5$ for
predominantly hydrodynamic operation, with the asperity-contact
regime conventionally identified as $\Lambda_{\min}<1$
\citep[\S~12-13]{shigley2020}. As with the assembly-tolerance
barrier, the indicator is not evaluated for the AP 1.8 head here
but its specification closes the gap identified by the
construct-validity benchmark and broadens the coverage of the
procedure beyond the spring--cam dynamic interaction.

\subsection{Subsequent analytical extensions}\label{sec:roadmap-extensions}

Subsequent extensions of the work will address (i)~a
second-generation Monte Carlo propagation that adds correlated
operational sources (thermal-soak transients, oil condition,
valve-seat wear) and joint barrier-failure probabilities; (ii)~a
coupled thermal--structural finite-element analysis of the cylinder
head to refine the differential hot-clearance estimate of
Section~\ref{sec:res-quant} and tighten the
$c_{\mathrm{cold,req}}\in [1.13,\,1.60]$~mm bracket; and (iii)~the
integration of the resulting indicator set as a constraint in the
multi-objective optimisation of the cylinder-head geometry. These
extensions are sequenced behind the test-rig programme of
Section~\ref{sec:roadmap-rig}: until the indicator magnitudes for
the AP 1.8 head are confirmed experimentally, the analytical
refinements rest on assumptions that the rig data alone can
substantiate.
